\section{Conclusion}
\label{sec:conclusion}

Part of what looks like capacity exhaustion in continual PEFT is not, and the two
parts call for two different fixes. On a $15$-task stream under one fixed LoRA
adapter, an additive per-label offset on the verbalizer logits --- zero rank, a
handful of scalars --- recovers accuracy that rank-allocation methods are built to
protect; that is an \emph{output-layer} constraint. What the offset cannot recover
is the representation the shared adapter overwrote as later tasks arrived; that is a
distinct \emph{representation-layer} constraint. This paper measures both and
proposes a single fixed-budget method that relieves both: reallocate the rank across
task families so each family adapter drifts only under its own family, and add the
per-scope offset on top.

\emph{The two moves compose, at a bit-identical budget.} Splitting the tasks into
$K{=}4$ semantic families with a rank-$2$ adapter each --- exactly the parameters of
one shared rank-$8$ adapter --- beats the shared adapter by $+8.00$~pp across three
seeds. The per-scope offset adds $+6.58$~pp on its own, and still adds $+5.18$~pp on
top of grouping, so the two constraints are independent and neither fix subsumes the
other. End-to-end, the combined method beats the standard single-adapter baseline by
$+11.76$~pp, positive on every seed and roughly nine times the same-configuration
noise floor (Section~\ref{sec:combined}). The gain is a reallocation, not extra
capacity: the trainable-parameter count is bit-for-bit identical across all four
cells of the $2\times2$.

\emph{The output component is real and its information content is specific.} Under
one globally shared offset the identification gap is $3.13$~pp; under an offset
shared only within a verbalizer scope it is $0.00$~pp, CI $[-1.71,+1.71]$. Scope ---
free at inference, readable from the prompt --- recovers $+1.51$~pp, CI
$[+0.43,+2.74]$. The gap is exactly zero when verbalizer sets are disjoint
(Proposition~\ref{prop:disjoint}), which we verify in every run and which explains
why the per-task-head literature cannot see this term at all.

\emph{The objective is the wrong lever for the representation constraint.} An
orthogonality penalty of O-LoRA's form, verbalizer-logit anchoring, and gauge-fixed
anchoring all left task-agnostic accuracy unimproved, two of them harmfully. This is
what motivated changing the \emph{allocation} rather than the objective, and the
allocation move is what delivered the representation-layer gain the penalties could
not.

\paragraph{What we report at its upper bound, honestly.}
Both moves are scored at their upper bound: the offset is fit against the current
model and family routing uses the scope at evaluation. The deployable offset sources
each carry a measured cost on this stream. Storing per-task optima and applying the
scope's Chebyshev centre --- $33$ floats, no task index --- scores $5.08$~pp
\emph{below} using no offset at all, because a value fit at a task's boundary goes
stale as a monolithic adapter keeps moving; singleton scopes locate the cause with
aggregation, routing, and gauge ruled out by construction ($-4.18$~pp, CI
$[-6.61,-1.76]$). Computing the offset from the query batch instead gains $+4.40$~pp
and then inverts to $-0.67$~pp under the class-mix control we pre-registered against
our own balanced splits. We report both plainly. Grouping is positioned to soften
the staleness cost --- a per-family adapter drifts less than a monolithic one --- and
we mark the de-staling of stored offsets under grouping, and a task-free family
router, as the measurements that would close the deployment gap. Neither is claimed
here; the allocation result does not depend on them.

\paragraph{The one task that loses, and what it points to.}
The grouped allocation is not uniformly better. The $14$-way DBpedia loses $5.1$~pp
because rank $2$ is genuinely too small for it --- the only task negative in all three
seeds --- and it is the same task whose offset cell stays competitive under the
shared adapter's larger rank. We report it rather than average it away, and it is the
concrete motivation for an adaptive per-family rank: a topic family with fourteen
labels should draw more of the budget than a binary one. That is the natural next
rung, and unlike the offset-deployment questions it is a pure allocation question the
same method already frames.

\paragraph{Scope.}
One stream (Order-4), one backbone (T5-large), one budget (LoRA $r{=}8$, $2.36$M
trainable scalars), three seeds, generative PEFT with a shared decoder and
restricted-argmax scoring. Family routing and the offset are scored at their upper
bound, as stated. Every comparison is internal: arms differ in one component at an
identical budget with identical splits and seeds. We ran no published continual-PEFT
baseline and make no claim against one, and the official test files were never
opened, so our absolute accuracies do not belong in leaderboard tables. The budgeted
quantization result is proved for binary shared verbalizers and remains open beyond
them.
